% ID: SEL-D0009-BB-E03
% ESTADO: cotejada. Segunda pasada visual completa y borrador compilado revisados.
% Fuente: Selectividad/matematicas_aplicadas/ORD_RESER_Modelo-3-Enunciados (Mat-Ap-2021).pdf
% SHA-256: 1dc0455361ce91a1157a7b3e4a06db95bc988e8dd0c9756e9aeb74b8cbb73a81
% Página: 1; bloque B; ejercicio original 3.

Se considera la función
\[
f(x)=\begin{cases}
-2x+2a & \text{si }-4\leq x\leq -2\\
-2x^2-4a & \text{si }-2<x\leq 2\\
-8x+b & \text{si }2<x\leq 3
\end{cases}
\]
\begin{enumerate}[label=\alph*)]
  \item \textbf{(1 punto)} Calcule los valores $a$ y $b$ para que la función
  sea continua en su dominio. Para esos valores, \textquestiondown{}es $f$
  derivable?
  \item \textbf{(0.8 puntos)} Para $a=-2$ y $b=16$, estudie la monotonía
  de la función $f$ y calcule sus extremos relativos y absolutos.
  \item \textbf{(0.7 puntos)} Para $a=-2$ y $b=16$, calcule el área del
  recinto limitado por la gráfica de $f$, el eje $OX$ y las rectas $x=-2$
  y $x=2$.
\end{enumerate}
